\documentclass[a4paper,12pt,twoside]{article}
\usepackage[utf8]{inputenc}

\usepackage{listings}
\usepackage{xcolor}

%New colors defined below
\definecolor{codegreen}{rgb}{0,0.6,0}
\definecolor{codegray}{rgb}{0.5,0.5,0.5}
\definecolor{codepurple}{rgb}{0.58,0,0.82}
\definecolor{backcolour}{rgb}{0.95,0.95,0.92}

%Code listing style named "mystyle"
\lstdefinestyle{mystyle}{
  backgroundcolor=\color{backcolour},   commentstyle=\color{codegreen},
  keywordstyle=\color{magenta},
  numberstyle=\tiny\color{codegray},
  stringstyle=\color{codepurple},
  basicstyle=\ttfamily\footnotesize,
  breakatwhitespace=false,
  breaklines=true,
  captionpos=b,
  keepspaces=true,
  numbers=left,
  numbersep=5pt,
  showspaces=false,
  showstringspaces=false,
  showtabs=false,
  tabsize=2
}

%"mystyle" code listing set
\lstset{style=mystyle}

% standard incantations
\usepackage[T1]{fontenc}
\usepackage[utf8]{inputenc}
\usepackage{lmodern}
\usepackage[german]{babel}
\usepackage{csquotes}

% clickable links in the PDF
\usepackage{xcolor}
\usepackage{hyperref}
\hypersetup{
	colorlinks=false,
	allbordercolors={white}
}
\usepackage{float}
% bibliography
\usepackage[sorting=none]{biblatex}
\addbibresource{literatur.bib}

% glossary
\usepackage[xindy]{glossaries}
\newglossaryentry{souveraenitaet}{%
  name={Souveränität},%
  description={Der Begriff Souveränität, auch „Staatshoheit“, wird im innerstaatlichen Recht und in der politischen Theorie verwendet, um die oberste Kompetenz zur Machtausübung im Innern eines Staates zu bezeichnen\cite{grundgesetzSouv}.}}

\newglossaryentry{kks}{%
  name={KKS},%
  description={künstliche Währung zum Ausgleich von Preis­niveau-Unterschieden zwischen den Mitgliedstaaten der Europäischen Union; ein Kaufkraftstandard (KKS) entspricht der durchschnittlichen Kaufkraft eines Euro in der Europäischen Union (EU-27)}}

\newacronym{eu}{EU}{Europäische Union}

\makeglossaries

% add literatur to toc
\usepackage[nottoc]{tocbibind}

% graphics and images
\usepackage{graphicx}
\usepackage{subfigure}
\usepackage{wrapfig}

% color packages
\usepackage{color, colortbl}
\definecolor{Gray}{RGB}{220,220,220}
\definecolor{EUBlue}{RGB}{45,172,227}
\definecolor{White}{RGB}{255,255,255}
\usepackage[first=0,last=9]{lcg}
\newcommand{\ra}{\rand0.\arabic{rand}}

% multirow table
\usepackage{multirow}

% footnote package
\usepackage{tablefootnote}

\lstset{
  captionpos=b,
  commentstyle=\color{UStuttDarkGreen},
  frame=single,	                   % adds a frame around the code
  keepspaces=true,
  %keywordstyle=\color{UStuttDarkBlue},
  showspaces=false,
  showstringspaces=false,          % underline spaces within strings only
  showtabs=false,
  stringstyle=\color{UStuttDarkBlue},
  tabsize=2
}

% --------------------------------------------------------------------
% Definitions of title informations
% --------------------------------------------------------------------
\newcommand{\HRule}[1]{\rule{\linewidth}{#1}}

\makeatletter
\def\printtitle{	
    {\centering \@title\par}}
\makeatother			

\makeatletter
\def\printauthor{
    {\centering \large \@author}}
\makeatother

% --------------------------------------------------------------------
% Config Title & Author
% --------------------------------------------------------------------
\title{
\HRule{0.5pt} \\
\LARGE \textbf{\uppercase{Graphauswertung}}
\HRule{2pt} \\ [0.5cm]
\normalsize \textsc{Ermittlung der Stimmungsmacher im Bundestag} \\
\normalsize \textsc{Markus Christopher Glutting, Marie Bittiehn, Miriam Lischke}
\\[2.0cm]
\normalsize \today
}

\author{
\normalsize Betreut von
\normalsize Prof. Dr. Thomas Hoppe\\
\normalsize Informationssysteme\\
\normalsize M.Sc. Angewandte Informatik\\
\normalsize Hochschule für Technik und Wirtschaft\\
\normalsize Treskowallee 8, 10318 Berlin, Deutschland\\
}

\begin{document}
% ------------------------------------------------------------------------------
% Maketitle
% ------------------------------------------------------------------------------
\thispagestyle{empty}
\printtitle
  	\vfill
\begin{figure}[H]
    \centering
    \includegraphics[width=200px, keepaspectratio]{logos/bundestag.png}
\end{figure}
  	\vfill
\printauthor		
\newpage

\pagenumbering{roman}
%\setcounter{page}{3}

\setcounter{tocdepth}{2}
\tableofcontents
\newpage

\listoffigures

\listoftables

\pagenumbering{arabic}
%\setcounter{page}{6}

\section{Einleitung}
% Was war die Aufgabe?

Der folgende Teilabschnitt der Ausarbeitung beschäftigt sich mit dem sechsten Teilprojekt: die Graphauswertung. Die Projektgruppe, welche an dem Teilprojekt gearbeitet hat, besteht aus Markus Glutting, Miriam Lischke und Marie Bittiehn.

\subsection{Hintergrund}
Das Teilprojekt ``Graphauswertung'' baut auf den Ergebnissen der Teilprojekte ``Interaktion zwischen Personen'' (Gruppe 4) und ``Interaktion zwischen Fraktionen'' (Gruppe 5) auf. Genauer formuliert, besteht die Aufgabenstellung darin, die Graphen, welche von Gruppe 4 und 5 erstellt werden, auszuwerten und die Ergebnisse der Auswertung der nachfolgenden Gruppe 8 für ihre Benutzeroberfläche zur Verfügung zu stellen.

\subsection{Problemstellung}
Durch die Auswertung der Graphen sollen die sogenannten ``Stimmungsmacher'' im Bundesstag ermittelt werden. Unter einem Stimmungsmacher ist im vorliegenden Kontext eine Person gemeint, welche viel mit vielen verschiedenen Personen redet und somit eine Stimmung verbreitet. Ob diese verbreitete Stimmung positiv oder negativ ist, ist dabei nicht entscheidend.

Neben der Ermittlung von Stimmungsmachern sollen ebenfalls simple mathematische Analysen auf den Graphen durchgeführt werden. Dadurch soll eine Gesamtbetrachtung der Sitzungen einer Wahlperiode ermöglicht werden, ebenso wie die Option die Sitzungen in Vergleich zueinander stellen zu können.

\subsection{Zielsetzung}
% Lernziele darstellen

Um Stimmungsmacher zu ermitteln, soll der PageRank-Algorithmus (siehe Kapitel \ref{sec:7-grundlagen}) verwendet werden. Dies bietet sich an, da Stimmungsmacher gleichbedeutend sind zu Personen, welche viele Nachrichten mit einem positiven und/oder einem negativen Sentiments empfangen bzw. versenden. Mithilfe des PageRank-Algorithmus werden eben diesen Personen bzw. Knoten im Graphen hohe Ränge vergeben, wodurch sie identifiziert werden können.

Für die mathematischen Analysen sollen Berechnungen auf den gesamten Graphen durchgeführt und statische Größen wie bspw. der Median oder ein Quartil der Sentiments berechnet werden.

Die erstellten Analysen sollen der nachfolgenden Projektgruppe (Nutzeroberfläche) über Schnittstellen zur Verfügung gestellt werden.

Für die Bearbeitung des Teilprojekts wurden zu Projektbeginn folgende Lernziele identifiziert:
\begin{itemize}
  \item Kennenlernen des PageRank Algorithmus
  \item Kennenlernen von graphenbasierten Datenbanken, insb. Neo4j mit Cypher
  \item Wissensaufbau im Bereich Backend Web Applikationen
\end{itemize}

\subsection{Prozess}
Die Bearbeitung des Teilprojekts lässt sich in vier Phasen gliedern:
\begin{itemize}
  \item Phase 1 (Oktober 2020): Thematische Einarbeitung und Projektplanung
  \item Phase 2 (November - Dezember 2020): Implementierung der geplanten Features
  \item Phase 3 (Januar 2021): Anpassungen in Rücksprache mit der nachfolgenden Projektgruppe (Nutzeroberfläche)
  \item Phase 4: (Februar 2021): Dokumentation
\end{itemize}


Die Umsetzung des Projekts erfolgte arbeitsteilig durch alle Mitglieder des Projektteams.
Darüber hinaus bestehen folgende Verantwortlichkeiten:


\begin{table}[ht]
\caption{Gruppe 7 (Graphauswertung) - Aufgabenverteilung}
\label{tab:zeittafelEU}
\begin{tabular}{|p{9cm}|p{5cm}|}
\hline
\rowcolor{Gray}
Aufgabe & Gruppenmitglieder \\
\hline
Server-Administration & Markus Glutting \\
Implementierung: PageRank & Marie Bittiehn \\
Implementierung: Stammdaten & alle \\
Implementierung: Statistische Analysen & Miriam Lischke \\
Organisation und Projekt Setup & Markus Glutting \\
Dokumentation & alle \\
\hline
\end{tabular}
\end{table}


Die Entwicklung erfolgte in wöchentlichen Zyklen.
Zu Beginn eines jeden Zyklus fand die wöchentliche Plenarsitzung im Rahmen der Lehrveranstaltung statt.
Im Anschluss daran erfolgte eine gruppeninterne Absprache, in der die zuletzt abgeschlossenen Aufgaben besprochen und die nächsten zu erledigenden Aufgaben identifiziert wurden.
Diese Aufgaben wurden mithilfe von GitHub Issues dokumentiert und der verantwortlichen Person zugewiesen.
Bis zur nächsten Plenarsitzung erfolgte die eigenverantwortliche Bearbeitung der Aufgaben.



\section{Grundlagen}
\label{sec:7-grundlagen}

In diesem Kapitel werden die notwendigen Kenntnisse für das Verständnis der nachfolgenden Kapitel vermittelt. Zuerst wird der PageRank-Algorithmus erläutert mit Fokus auf eine effizientere Berechnung der Ränge mithilfe von Matrizenmultiplikationen. Nachfolgend wird auf die wichtigsten verwendeten Entwicklungsframeworks, -Tools und -Konzepte eingegangen.

\subsection{PageRank}
Der PageRank-Algorithmus ist ein Algorithmus zur Gewichtung von Knoten innerhalb eines Netzwerks anhand der Anzahl ihrer eingehenden Beziehungen. Dabei geht in die Berechnung des PageRanks eines Knoten, neben der Anzahl an eingehenden Beziehungen, ebenso der PageRank der auf ihn verweisenden Knoten mit ein. 

Aufgrund der daraus entstehenden direkten Abhängigkeit des PageRank eines Knoten von den PageRanks der auf ihn verweisenden Knoten, ist vor allem bei großen Netzwerken ein genaue Berechnung aller PageRanks nicht immer in absehbarer Zeit möglich. Aus diesem Grund werden die Werte für die PageRanks bei der Berechnung meist iterativ angenähert. Dabei wird allen Knoten ein einheitlicher Startwert als PageRank vergeben. Meist wird als Startwert der Wert $1/N$ verwendet, wobei $N$ die Anzahl aller Knoten des Netzwerks ist. Danach wird der PageRank für alle Knoten mehrfach berechnet und so die Werte iterativ angenähert.

Eine Möglichkeit dieser iterativen Annäherung ist ein mehrfaches Iterieren über das gesamte Netzwerk und die rekursive Berechnung des PageRanks für jeden Knoten. Da im vorliegenden Anwendungskontext mit mehreren hundert Knoten und zig Tausend von Beziehungen zwischen den Knoten gerechnet werden kann, wäre das Iterieren über alle Knoten und Beziehungen nicht sehr effizient. Stattdessen können allerdings Matrizenmultiplikationen und eine Eigenvektorberechnung für die PageRank-Berechnung verwendet werden.

Das Netzwerk, welches aus mathematischer Sicht als gerichteter Graph betrachtet werden kann, kann als eine quadratische, stochastische Matrix abgebildet werden \cite{pagerank_eigenvector}. Stochastisch bedeutet in diesem Kontext, dass die Summen der Spalten der Matrix alle 1 betragen \cite{pagerank_eigenvector}. Wenn im Graphen von einem Knoten $j$ eine Beziehung  zum Knoten $i$ besteht, dann wird in der Matrix an der Stelle $ij$ der Wert $1/d^{+}(j)$ eingetragen, wobei $d^{+}(j)$ für den Ausgangsgrad von $j$ steht \cite{pagerank_eigenvector}. Wenn keine Beziehung besteht, wird an der Stelle $ij$ eine 0 eingetragen \cite{pagerank_eigenvector}. Für den Fall, dass ein Knoten gar keine ausgehenden Beziehungen besitzt, wird in seiner gesamten Spalte in der Matrix der Wert $1/N$ eingetragen \cite{pagerank_eigenvector}. Dadurch wird verhindert, dass der PageRank in Sackgassen-Knoten gewissermaßen ``versickert''. 

Wurde der Graph nun als Matrix abgebildet, so kann der PageRank iterativ mithilfe von Matrixmultiplikationen berechnet werden. Dabei wird der PageRank als Vektor der Länge $N$ dargestellt, meist befüllt mit dem Startwert $1/N$ für alle Knoten. Statt einer gesamten Iteration über das Netzwerk und den rekursiven Berechnungen neuer PageRanks für alle Knoten, muss nun nur eine simple Multiplikation der Matrix mit dem Vektor durchgeführt werden. Das Ergebnis der Multiplikation ist wieder ein Vektor und dieser enthält die neuen PageRanks aller Knoten. Nach mehrfacher Multiplikationen der Matrix mit dem PageRank-Vektoren ändern sich die Werte innerhalb des Vektor nicht mehr. Es wurde iterativ ein dominanter Eigenvektor der Matrix berechnet, welcher gleichzeitig den endgültigen PageRanks der Knoten des Netzwerks entspricht \cite{pagerank_eigenvector}. 

Auf gleiche Art und Weise kann der sogenannte ``Reverse PageRank'' berechnet werden. Bei dem Reverse PageRank handelt es sich um den PageRank berechnet auf dem Reverse-Graphen des Originalgraphen \cite{reverse_pagerank}. Dies bedeutet, dass der Reverse PageRank nicht basierend auf den eingehenden Beziehungen der Knoten berechnet wird, sondern auf den ausgehenden Beziehungen. 

Um im ersten Schritt den Reverse-Graphen des Originalgraphen zu erhalten, werden die Richtungen der Beziehungen zwischen den Knoten des Originalgraphen umgedreht \cite{reverse_pagerank}. Danach wird der so erzeugte Reverse-Graph ebenfalls als Matrix abgebildet und die PageRank-Werte der Knoten werden mithilfe von Matrizenmultiplikationen berechnet. Die dadurch ermittelten Werte für den Reverse PageRank ähneln den Ergebnissen des klassischen PageRank dahingehend, dass die einzelnen Werte gleich sind, aber anders verteilt. Dies ist dadurch zu erklären, dass bei der Berechnung des Reverse PageRank zwar die Richtungen der Beziehungen umgedreht werden, jedoch keine Beziehungen entfernt oder hinzugefügt werden. Aus diesem Grund sind die Werte gleich, aber durch das Umdrehen der Beziehungen sind sie anders verteilt.

\subsection{Entwicklungsframeworks, -Tools und -Konzepte}

Vor der Wahl eines Entwicklungsframeworks galt es zunächst zu entscheiden, welche Programmiersprache verwendet. Innerhalb der Gruppe war die Python-Bibliothek numpy bekannt, welche gut zur performanten Berechnung von Matrizen geeignet ist.
Daher fiel die Entscheidung auf die Programmiersprache Python. Zur Bereitstellung von HTTP-Schnittstellen für die Benutzeroberfläche wurde das Microframework Flask verwendet.
Flask zeichnet sich dadurch aus, dass es einen sehr schlanken Kern hat, der nach Bedarf erweitert werden kann. Dadurch unterscheidet es sich stark von anderen Frameworks wie bspw. Django, welches stattdessen eher darauf abzielt, ein vollständiges Framework für Webanwendungen zu sein.

Der Zugriff auf die Neo4j Datenbaken erfolgt mithilfe des Neo4j Python Treibers, welcher im Python Paket neo4j bereitgestellt wird. Es wurde kein Object Graph Mapper (z.B. neomodel) verwendet, da eines der Lernziele darin bestand, die Abfragesprache Cypher kennenzulernen.



\section{Anforderungsanalyse und Konzept}
% Wie sind Sie vorgegangen?

Zur Identifikation der Anforderungen

\subsection{Architektur}

TODO: Diagramm + Beschreibung

\subsection{Schnittstellen}

1) HTTP Endpoints für Nutzeroberfläche

2) Datenbankzugriff auf Neo4j Datenbanken

\section{Implementierung und Bereistellung}

\subsection{Implementierung}
Die Backend-Anwendung wurde als Python Flask Applikation entwickelt. Der gesamte Quellcode ist im GitHub Repository \textit{Sentiments-of-Bundestag/ Graphenauswertung}~\cite{github_graphenauswertung} verfügbar.

Die Applikation wird durch Ausführen der Datei \lstinline{app.py} gestartet.
Diese stellt zunächst die Verbindung mit den beiden Neo4j Datenbanken her.
Zur Kommunikation mit den beiden Datenbanken existiert jeweils eine Klasse, die den Zugriff auf die jeweilige Datenbank verwaltet und Methoden zur Abfrage der Inhalte der Datenbank bereitstellt.
Die Zugangsdaten für die Datenbanken werden als Umgebungsvariablen bereitgestellt. Diese werden beim Start der Applikation aus der Datei \lstinline{.env} eingelesen.
Aus Sicherheitsgründen, ist diese Datei nicht Teil der Quellcodeverwaltung.
Einige komplexere Prozeduren wie bspw. die Berechnug des PageRank wurden ausgelagert und werden von den Datenbankklassen importiert.
Die Defintion der Schnittstellen erfolgt in Form von Methoden, die mit dem Pfad annotiert sind, unter dem die jeweilige Schnittstelle erreichbar ist.
Aufgrund der großen Datenmenge und der damit verbundenen hohen Laufzeit werden alle Anfragen für 12 Stunden gecached.
Dazu wird die Bibliothek Flask-Caching mit dem Cache Typ SimpleCache verwendet, welcher einem In-Memory Python Dictionary entspricht.\cite{flask_caching}
Die Implementierung der Schnittstellenmethoden besteht in der Regel aus einem Aufruf einer Methode der Datenbankklassen, wobei das Ergebnis zur Übertragung in der HTTP Response in das JSON Format umgewandelt wird.

Die Methoden der Datenbankklassen sind alle nach dem gleichen Muster aufgebaut.
Die nachfolgende Methode, welche die verfügbaren Fraktionen abfragt, kann als Beispiel dienen.

\begin{lstlisting}
def get_factions(self,
                 sentiment_type="NEUTRAL",
                 session_id=None):

    where = self.get_where_clause(sentiment_type, session_id)
    with self.driver.session() as session:
        query = \
            "MATCH (sender:{0})-[r:{1}]-(recipient:{0}) " \
            "{2} " \
            "with sender, " \
            "collect(distinct r.sessionId) as sesList, " \
            "collect(r) as rlist unwind rlist as r " \
            "RETURN DISTINCT sender.name as name, " \
            "sender.size as size, " \
            "sender.factionId as factionId, " \
            "sesList as sessionIds" \
            .format(NODE_FACTION, REL_COMMENTED, where)

    factions = session.run(query)
    return factions.data()
\end{lstlisting}

Die Methoden sind stets so aufgebaut, dass zunächst eine Datenbanksitzung geöffnet wird.
Zur Datenbankabfrage wird eine Cypher Query konstruiert. Diese hängt meist von den Parametern der HTTP-Anfrage ab, welche in der \lstinline{WHERE} Klausel der Query abgebildet werden.
Da viele Methoden die gleichen Parameter unterstützen, ist die Konstruktion der \lstinline{WHERE} Klausel in eine eigene Methode ausgelagert.
Die Query wird anschließend an die jeweilige Neo4j Datenbank geschickt. Das Resultat wird in einigen Fällen noch weiter verarbeitet und dann als Resultat zurückgegeben.
Mit dem Abschluss der Methode wird auch die Datenbanksitzung beendet.

\subsection{Bereistellung}
Zur Bereitstellung der Anwendung auf einem Server wurde ein Docker Image erstellt, mithilfe dessen die Anwendung gestartet werden kann.
Dabei wurde wie in \cite{flask_docker} beschrieben vorgegangen.
Das Docker Image basiert auf dem Standard Docker Image für Python und wird gebaut, indem zunächst die Bibliotheken, von denen die Anwendung abhängig ist installiert werden.
Anschließend wird der Quellcode und die Konfigurationsdateien in das Docker Image kopiert.
Zum Start der Anwendung wird ein Shell Skript aufgerufen, welches den uWSGI Applikationsserver startet.
HTTP-Anfragen auf dem Standard Port 80 werden mithilfe von nginx an den uWSGI Server weitergeleitet.
Um den Start der Anwendung ohne Angabe des Port Mappings zu gestalten wurde außerdem eine \lstinline{docker-compose} Datei erstellt.

Die Bereitstellung auf dem HTW Server erfolgt durch folgende Schritte:
\begin{itemize}
  \item Zugriff auf den Server via \lstinline{ssh}
  \item Herunterladen der neuesten Version des Quellcodes mithilfe von \lstinline{git pull}
  \item Erstellung des neuen Docker Images mithilfe von \lstinline{docker-compose build}
  \item Start der Applikation mithilfe von \lstinline{docker-compose up}
\end{itemize}

Die zuvor beschriebenen Schnittstellen wurden in der \lstinline{README} des Repositories beschrieben und so der nachfolgenden Gruppe (Nutzeroberfläche) zur Verfügung gestellt.
Die Nutzeroberfläche greift auf die Schnittstellen zu, um die Analysen auf der Website verfügbar zu machen.

\section{Fazit}

\subsection{Zusammenfassung}
% Was wurde erreicht?

\subsection{Lernziele}
% Lernziele auswerten (erreicht: ja/nein)

\subsection{Ausblick}
% Was blieb offen?
% Wo liegen Verbesserungsmöglichkeiten?

\pagenumbering{roman}
%\setcounter{page}{6}

\printbibliography[title={Literaturverzeichnis}]
\addcontentsline{toc}{section}{Literaturverzeichnis}
\newpage

\printglossaries
\addcontentsline{toc}{section}{Glossar}
\newpage

\end{document}